% =====================================================================
% doc-example.tex — sheffielddoc.sty demonstration handout
% Compile:  pdflatex doc-example.tex  (twice, for TOC/refs)
% Needs:    sheffielddoc.sty + assets/crest-white.png beside this file
% =====================================================================
\documentclass[10pt,a4paper]{article}
\usepackage[most]{tcolorbox}   % load before the theme
\usepackage{longtable,multirow,microtype}
\usepackage[hidelinks]{hyperref}
\usepackage[watermark=arrows]{sheffielddoc}

\sheffdocheader{Introduction to Aerodynamic CFD}

\begin{document}

% --- Title page ------------------------------------------------------
\ShefDocTitle
  {Introduction to Aerodynamic CFD}
  {A rocket fin, from STEP file to lift coefficient}
  {Student handout \quad$\cdot$\quad ANSYS 2023 R2 / 2024 R1 \quad$\cdot$\quad a three-hour workshop}

\ShefDocIntro
  {What this document is}
  {A complete, click-by-click record of the workshop: every number you type,
   and the reason it is that number rather than another.}
  {What you need}
  {ANSYS Workbench with Fluent and Discovery, the file \file{fin.step} from
   the course folder, and roughly 6\,GB of free disk.}

\vfill
{\footnotesize\color{MidGrey}
 Version 1.0. Prepared for undergraduate teaching.\par}
\clearpage

% --- Front matter ----------------------------------------------------
\section*{How to use this handout}
\markboth{How to use this handout}{}
\addcontentsline{toc}{section}{How to use this handout}

The workshop moves through four stages --- geometry, mesh, solution,
post-processing. Each stage has the same shape:

\begin{itemize}
  \item \ui{Bold sans} is a button, tab, panel or dropdown you can click.
  \item \val{Purple bold} is a value you type or select.
  \item \menu{File \sep Preferences} is a menu path: click \ui{File}, then
    \ui{Preferences}.
  \item \file{monospace} is a file name or a name in a tree.
\end{itemize}

\begin{center}
\small
\begin{tabular}{@{}L{0.17\textwidth}L{0.74\textwidth}@{}}
\toprule
\textbf{Element} & \textbf{What it does}\\
\midrule
\textbf{Why} & the theory that makes the next click defensible.\\
\textbf{Do this} & the numbered sequence of clicks and the exact values.\\
\textbf{Checkpoint} & what the screen should look like when the step worked.\\
\textbf{If it goes wrong} & the failures that actually happen, and the fix.\\
\bottomrule
\end{tabular}
\end{center}

\clearpage
\tableofcontents
\clearpage

% --- Body ------------------------------------------------------------
\section{The case: a fin, and the four numbers that define it}

\subsection{Geometry}

The course folder contains \file{fin.step}: a model rocket fin with a swept
leading edge, a trapezoidal planform and a sharp trailing edge.

\begin{why}
\boxlab{Why the Reynolds number comes first}
Every decision downstream --- the size of the first cell at the wall, the
density of the wake --- depends on the Reynolds number. It is fixed by the
chord, the speed and the air, and nothing later in the workflow can repair
choosing it carelessly.
\end{why}

\ShefStep{Set the operating point}

Run the case at \val{U = 40 m/s}, sea level. With $c = 250$\,mm this gives
$Re \approx 5.1\times10^{5}$ --- fully turbulent over most of the fin.

\begin{vals}
\boxlab{Values to enter}
\begin{itemize}
  \item Velocity magnitude: \val{40 m/s}
  \item Density (constant): \val{1.225 kg/m\textsuperscript{3}}
  \item Viscosity: \val{1.7894e-05 kg/m\,s}
\end{itemize}
\end{vals}

\begin{chk}
\boxlab{Checkpoint}
The boundary-condition panel reads \ui{velocity-inlet}, magnitude
\val{40}, and the residual monitor starts near $10^{-1}$.
\end{chk}

\begin{fail}
\boxlab{If it goes wrong: the mesh will not fill}
If the fill fails with a negative-volume error, the fluid body is not
watertight: go back to \S2.4, run the fault-isolating tool, and re-check
the share topology before retrying.
\end{fail}

\begin{asid}
{\small\color{MidGrey}Aside: the fin is thin, so the spanwise flow is
weak; a symmetry treatment would be defensible but we model the full
span for cleaner post-processing.}
\end{asid}

\subsection{What the solver actually solves}

For an incompressible, steady case Fluent solves the Reynolds-averaged
momentum equations plus continuity; see the handout appendix for the
discretisation settings and their defaults.

\section{Closing}

\ShefDocContact{mech-eng-workshops@sheffield.ac.uk}

\end{document}
